% Paper fragment for NAV.cls. Requires amsmath and booktabs; NAV.cls supplies natbib.
\section{Pl\'eiades imagery: morphology and conditional source-location analysis}
\label{sec:pleiades-conditional}

On 23 March 2014, Pl\'eiades 1A acquired four 0.5~m southern
Indian Ocean scenes. At the ATSB's request, Geoscience Australia analysed
imagery from French military intelligence and CNES~\citep{minchin2017pleiades}.
Twelve of at least 70 objects were rated 5 (probably not natural objects or
features). This did not identify aircraft debris. Targets 06, 18, 19, 26 and
27 were retained because repeated public-panel segmentation was stable (15/18
size-qualified trials for 06; 18/18 otherwise), not because they resembled aircraft.

Eight equal-sized families covered 777 parts, matched random geometry, ocean
controls, and whole, documented or random fuselage sections. Boeing
geometry~\citep{boeing2024acap} supplied shapes; controls included containers,
nets, cargo/plastic/timber, hulls and fish-aggregating
devices~\citep{hapag2016container,iotc2023observer,lebreton2018plastic}.
\citet{murphy2013structures} gives six reproducible section boundaries; they

Each family contained 3,372 masks (26,976; 134,880 family--object
comparisons). Fuselage orientation was isotropic; 20\% were fully visible and
the remainder used 15--100\% linear-waterline coverage. Eligibility required
area, length and width within both limits. These were sensitivities, not
probabilities. Principal-axis-normalised masks used four reflections. For
candidate $c$, target $o$, family $g$, masks $T_o,P_o$, reflections
$\mathcal F$, and eligible set $E_o$,
\begin{equation}
\begin{aligned}
  \operatorname{IoU}(A,B)
    &=\frac{|A\cap B|}{|A\cup B|},\\
  \ell(c,o)
    &=1-\frac{
      \max_{f\in\mathcal F}\operatorname{IoU}(T_o,f(c))+
      \max_{f\in\mathcal F}\operatorname{IoU}(P_o,f(c))
    }{2},\\
  \ell^{*}_{g,o}
    &=\min_{c\in g\cap E_o}\ell(c,o).
\end{aligned}
\label{eq:pleiades-mask-loss}
\end{equation}
Lower loss is greater overlap, not an identity probability.

Extended minima were an ocean net/rope control (06), matched random geometry
(18/19), a random-pose wing (26), and documented section 47 (27). Object-06
ocean and random losses differed by 0.0007. Section 47 scored 0.3058, but
0.3556 when fully visible. The whole shell ranked lowest for none.
Dimension-matched random controls therefore left 777 shapes admissible but
non-specific.

Object 18's complete-wing PCA minimum was unconstrained. In a post-hoc
hydrostatic sensitivity, none of 168 original flooding states matched; 10/24
multistart refinements met the 25\%-per-slope and strict-size criteria. All
required an absent engine, flooded midspan/tip cells and inboard-only buoyancy.
The closest four-height RMSE was 0.796~m and PCA loss 0.3421, versus 0.2730
unconstrained. The cells are not Boeing tank bays. Integral tanks and protected
engine separation make residual buoyancy possible, but atmospheric venting and
no documented sealed inboard bay leave it
unsupported~\citep{aaib2010ba38,motmalaysia2018,ntsb2014fueltank}. Target
fitting and unequal budget exclude it from the ranking.

\begin{table}[t]
  \centering
  \footnotesize
  \caption{Equal-budget Pl\'eiades morphology controls. Lower loss is better.}
  \label{tab:pleiades-morphology-controls}
  \begin{tabular}{lrrr}
    \toprule
    Family & Eligible objects & Mean lowest loss\textsuperscript{a} &
      Objects at lowest loss \\
    \midrule
    Flooding-model wing            & 1/5 & 0.3286 & 0 \\
    Random-pose wing               & 5/5 & 0.3232 & 1 \\
    Other shortlisted 777 parts    & 4/5 & 0.3704 & 0 \\
    Matched random geometry        & 5/5 & 0.3267 & 2 \\
    Generic ocean objects          & 4/5 & 0.3532 & 1 \\
    Whole 777 fuselage               & 4/5 & 0.3472 & 0 \\
    Documented 777 fuselage sections & 5/5 & 0.3318 & 1 \\
    Random 777 fuselage sections     & 4/5 & 0.3388 & 0 \\
    \bottomrule
  \end{tabular}

  \vspace{2pt}
  \parbox{0.94\linewidth}{\footnotesize\textsuperscript{a}Mean over eligible
  objects only; family means are therefore not directly comparable.}
\end{table}

For source location, \citet{griffin2017drift} propagated debris forward
conditional on some objects being from 9M-MRO. A clean-sheet ensemble released
particles from 5,289 cells within $\pm100$~NM of the seventh arc. BRAN2016,
OSCAR v2~\citep{dohan2021oscar}, and GLORYS12 currents plus WAVERYS Stokes
drift~\citep{copernicus2023glorys12,copernicus2024waverys} were alternatives.
Common NCEP winds~\citep{kalnay1996reanalysis}, three windages and 5~NM/day
dispersion were used. All twelve location likelihoods were averaged rather than
drawing one object; the five retained targets formed a sensitivity.

Modes were 35.418$^{\circ}$S, 92.886$^{\circ}$E for BRAN
(24,309~km$^2$ 90\% region); 34.917$^{\circ}$S, 92.047$^{\circ}$E for OSCAR
(18,393~km$^2$); and 35.391$^{\circ}$S, 92.163$^{\circ}$E for
GLORYS12/WAVERYS (22,123~km$^2$). Pairwise total variation was 0.484--0.570 and
mode separation 53.8--94.4~km. An equal-prior mixture has mode
35.218$^{\circ}$S, 92.247$^{\circ}$E and a 28,768~km$^2$ region. It is
neither independent-evidence multiplication nor a calibrated posterior. The
PDFs enter the integrated estimator only under the conditional image-debris
hypothesis and provide no unconditional update.
